\documentclass[11pt]{article}

\usepackage[a4paper,margin=25mm]{geometry}
\usepackage{amsmath,amssymb}
\usepackage{booktabs,longtable,array}
\usepackage[table]{xcolor}
\usepackage{microtype}
\usepackage[hidelinks]{hyperref}
\usepackage{enumitem}

\definecolor{ok}{RGB}{20,110,55}
\definecolor{warn}{RGB}{166,91,0}
\definecolor{lightgray}{RGB}{245,246,248}
\newcommand{\R}{\mathbb R}
\newcommand{\norm}[1]{\lVert #1\rVert_2}
\newcommand{\abs}[1]{\lvert #1\rvert}
\newcommand{\spec}{\operatorname{spec}}
\newcommand{\dist}{\operatorname{dist}}
\newcommand{\eps}{\varepsilon}
\makeatletter
\g@addto@macro\UrlBreaks{\do\.\do\_}
\makeatother
\newcommand{\lean}[1]{\nolinkurl{#1}}
\newcommand{\match}{\textcolor{ok}{\textbf{Match}}}
\newcommand{\qualified}{\textcolor{warn}{\textbf{Qualified}}}
\setlength{\emergencystretch}{3em}
\setlist{nosep,leftmargin=2em}

\hypersetup{
  pdftitle={Audit of body.tex against the Lean verification},
  pdfauthor={Codex},
  pdfsubject={Statement correspondence audit for FinitePrecisionLanczos}
}

\title{Audit of \texttt{body.tex} against the Lean verification}
\author{Statement-correspondence report}
\date{20 August 2026}

\begin{document}
\maketitle

\section*{Executive conclusion}

The principal mathematical claims agree with the Lean development, subject to
the formal scope described below.  All thirteen labeled theorem or lemma
statements in \texttt{body.tex} have corresponding proved Lean declarations;
no contradictory bound or missing hypothesis was found in those displayed
results.  The asymptotic notation $\lesssim_k$ is justified by stronger,
explicit polynomial bounds in the code.

The body as a whole should nevertheless not be described as literally proved
by Lean without qualification.  There are two material wording issues:

\begin{enumerate}
  \item The Lean development starts from an abstract family of real vectors,
  scalars, and error terms satisfying the \lean{LanczosRun} interface.  It does
  not prove that executing Algorithm~1 in IEEE floating point produces such an
  interface.
  \item Exact preservation of starting spectral weights holds between the
  original matrix $A$ and the unperturbed dilation $S$.  For the nearby exact
  Lanczos problem $\widetilde T$, Lean proves perturbative Wasserstein and
  separated-cluster mass bounds, not exact equality of starting weights.
\end{enumerate}

Thus the appropriate verdict is: \textbf{the labeled mathematical results
match; the algorithmic and spectral-weight prose needs qualification.}

\section*{Verification performed}

The following checks were run directly on the supplied sources.

\begin{itemize}
  \item \texttt{lake\ build} completed successfully (3,238 jobs).  The output
  contained linter warnings but no build error.
  \item \texttt{lake\ env\ lean\ Audit.lean} checked the principal declarations.
  Each depends only on \lean{propext}, \lean{Classical.choice}, and
  \lean{Quot.sound}, the standard logical axioms used by Lean/Mathlib.
  \item A source search found no \lean{sorry}, \lean{admit}, or user-declared
  \lean{axiom} in the Lean files.
  \item The pinned toolchain is Lean~4.32.2.
\end{itemize}

The authoritative formal input is the structure
\lean{FinitePrecisionLanczos.LanczosRun} in
\lean{Lanczos/Run.lean}.  Its important global hypotheses are
\[
  H^T=H,\qquad \eps\ge 0,\qquad k>0,\qquad k\eps\le \frac1{100},
\]
together with the five local arithmetic identities and error bounds,
near-normalization, nonnegative $\beta_j$, and strict positivity of the
internal subdiagonals $\beta_1,\ldots,\beta_{k-1}$.

\section*{Correspondence of labeled results}

In the table, ``match'' means that the body statement follows immediately
from the cited Lean declaration after replacing the explicit polynomial bound
by $\lesssim_k$.  Line numbers refer to the supplied \texttt{body.tex}.

\renewcommand{\arraystretch}{1.16}
\begin{longtable}{@{}>{\raggedright\arraybackslash}p{0.22\textwidth}
                    >{\raggedright\arraybackslash}p{0.32\textwidth}
                    >{\raggedright\arraybackslash}p{0.39\textwidth}@{}}
\toprule
Body result & Principal Lean declaration & Finding \\
\midrule
\endfirsthead
\toprule
Body result & Principal Lean declaration & Finding \\
\midrule
\endhead
\bottomrule
\endfoot

Local error bounds (lines 265--275)
& \lean{LanczosRun.local\_errors}
& \match.  Lean proves
$\norm{f_j}\le700j\eps\norm H$, $\abs{g_j}\le\eps$, and
$\abs{p_j}\le1200j\eps\norm H$.  The exact identity used for $p_j$ omits
$e^u_j$, confirming the stated cancellation. \\

A priori coefficient bounds (282--290)
& \lean{LanczosRun.coefficients}
& \match.  Each of $\norm{u_j},\norm{w_j},\abs{\alpha_j},\norm{z_j}$ and
$\beta_j$ is at most $100j\norm H$. \\

Gram commutator (447--464)
& \lean{LanczosRun.gram\_commutator\_conclusion}
& \match.  Lean proves both exact commutator identities, the support statement
for the diagonal-defect commutator, and explicit entry bounds. \\

Paige loss bound (521--529)
& \lean{LanczosRun.paige\_bound}
& \match.  Lean proves the absolute boundary product is at most
$6000k^3\eps\norm H$, which is exactly the body quantity after regrouping
$\beta_k$. \\

Paige descent (577--586)
& \lean{LanczosRun.paige\_descent}
& \match.  The two conclusions are both bounded by
$8t(6000k^3\eps\norm H)$ for a smaller prefix. \\

Ritz containment (709--716)
& \lean{LanczosRun.containment\_prefix},
\lean{LanczosRun.containment\_eigenvalue}
& \match.  Lean proves the stronger explicit radius
$[(8t^2+5)6000k^3+2100k^2]\eps\norm H$ for every prefix $T_t$. \\

Stabilized Ritz localization (784--793)
& \lean{LanczosRun.stabilized},
\lean{LanczosRun.terminal\_localization\_eigenvalue}
& \match.  The code gives
$3\beta_k\abs{y_k}$ plus an explicit $O(k^5\eps\norm H)$ term, and the
terminal conclusion when $\beta_k=0$. \\

Normalized local errors (850--859)
& \lean{LanczosRun.normalized\_local\_errors},
\lean{LanczosRun.normalized\_governing}
& \match.  The residual is at most $2600k^4\eps\norm H$ and both kinds of
weighted adjacent product are at most $4800k\eps\norm H$. \\

Normalized commutator (892--904)
& \lean{LanczosRun.normalized\_gram\_commutator}
& \match.  Upper triangularity and the exact skew identity are proved, with
norm bound $50000k^5\eps\norm H$. \\

Nilpotent isometry (950--961)
& \lean{LanczosRun.nilpotent\_isometry}
& \match.  Strict upper triangularity, $K^k=0$, $\norm K\le1$,
$\norm C\le2$, and all three displayed identities are formal theorems. \\

Physical dilation (1001--1017)
& \lean{LanczosRun.physical\_dilation},
\lean{LanczosRun.block\_recurrence}
& \match.  Lean proves the orthogonality, first-column, recurrence, and norm
claims, with $\norm{\widehat E}\le10^6k^8\eps\norm H$. \\

Symmetric completion (1110--1120)
& \lean{LanczosRun.greenbaum\_exact\_completion},
\lean{GreenbaumModel.prescribed\_subdiagonal\_positive}
& \match.  The leading block is exactly $T_k$, the completed matrix is
symmetric tridiagonal, and the operator difference is at most
$3\norm{\widehat E}$.  Positivity is guaranteed through the leading block. \\

Greenbaum backward model (1146--1159)
& \lean{LanczosRun.exists\_greenbaumModel},
\lean{GreenbaumModel.exact\_lanczos\_returns\_Tmat}
& \qualified.  The matrix, dimension, leading block, and
$3{,}000{,}000k^8\eps\norm H$ bound match.  Exact reproduction is proved for
the formal \lean{ExactLanczos} interface, not by verification of a particular
software implementation. \\

\end{longtable}

\section*{Scope difference: Algorithm~1 versus \texttt{LanczosRun}}

The body moves from kernel assumptions and Algorithm~1 to a completed run.
The Lean code does not contain floating-point values, a rounding operator
$\operatorname{fl}$, IEEE~754 semantics, or an implementation of Algorithm~1.
Instead, every run-level theorem accepts the local identities and bounds as
fields of \lean{LanczosRun}.  Consequently, statements such as

\begin{quote}
``Suppose Algorithm~1 completes $k$ iterations under the assumptions of the
model''
\end{quote}

are supported only after an external proof supplies the bridge
\[
 \text{execution of Algorithm 1}
 \quad\Longrightarrow\quad
 \text{an instance of \lean{LanczosRun}}.
\]
No such bridge is present in the supplied Lean project.  This does not weaken
the conditional matrix analysis; it limits what has been verified about an
actual floating-point program.

There are also two benign differences showing that the formal interface is
slightly more general than the paper's operational description.  Lean does
not assume $k\le n$, and when $\beta_k=0$ it does not require separately that
$z_k=\Delta_k=0$; it needs only the stated local recurrence identity and
error bound.  Every body run satisfying its conventions is still covered.

\section*{Material wording issue: starting spectral weights}

The prose at lines 819 and 1261 says that the construction or model
``preserves the starting spectral weights exactly.''  This is correct only
for the unperturbed dilation $S$, which is orthogonally similar to $k$ copies
of $A$.  Lean proves
\[
  \mu_{S,e_1}=\mu_{A,\bar v_1}
\]
in \lean{GreenbaumModel.starting\_spectralMass} (and proves equality of every
starting moment in \lean{starting\_moment}).

The nearby problem on which exact Lanczos reproduces $T_k$ is instead
$\widetilde T$.  Between $A$ and $\widetilde T$, Lean proves
\[
 \mathcal W^{\rm dual}_1
   (\mu_{A,\bar v_1},\mu_{\widetilde T,e_1})
 \le \sqrt{nk}\,\delta,
 \qquad
 \delta=\norm{S-\widetilde T}
 \le 3{,}000{,}000k^8\eps\norm H,
\]
not exact equality.  It also proves a cluster-mass estimate only when the
chosen physical cluster is separated by a gap $g>2\delta$:
\[
 \abs{\text{physical mass}-\text{model cluster mass}}
 \le \frac{\sqrt{nk}\,\delta}{g-2\delta}.
\]

A precise replacement for the disputed wording is:

\begin{quote}
``The unperturbed dilation preserves the starting spectral measure exactly;
the nearby completed problem differs from it by one operator perturbation,
which yields dual-Wasserstein and separated-cluster mass bounds.''
\end{quote}

With this change, the spectral interpretation matches the formal results.

\section*{Claims outside the Lean verification}

The following parts of the body are not contradicted by the code, but they
are not established by its declarations and should not be presented as
machine-checked consequences:

\begin{itemize}
  \item the IEEE single-precision experiments and all numerical values in the
  three figure discussions;
  \item historical and literature-priority claims;
  \item broad empirical statements about practical convergence and accuracy
  of Lanczos-based matrix-function approximations;
  \item the informal eigenvector-accuracy remark following the Ritz residual
  bound (standard perturbation reasoning, but not packaged in this Lean
  development).
\end{itemize}

The Lean development actually proves two useful spectral consequences that
the main backward theorem in the body does not state explicitly: every
eigenvalue of $\widetilde T$ lies within $\delta$ of $\spec(H)$, and the
starting measures satisfy the dual-Wasserstein and separated-cluster bounds
displayed above.

\section*{Recommended minimal edits to \texttt{body.tex}}

\begin{enumerate}
  \item After the sentence about the companion Lean proof (line 93), add:
  ``The formalization begins from the local real-arithmetic run interface
  displayed below; it does not formalize a floating-point implementation of
  Algorithm~1.''
  \item At lines 819 and 1261, replace unqualified exact preservation of
  ``the model'' weights by the precise $S$ versus $\widetilde T$ statement
  given above.
  \item In the consequence of Theorem~\texttt{thm:greenbaum}, describe exact
  Lanczos as the formal exact three-term-recurrence interface, or define that
  interface explicitly.
\end{enumerate}

After those edits, the paper-facing theorem statements and the formal scope
would be aligned without a material caveat.

\end{document}
